\chapter{Toy examples}
\label{ch:toy}
The approach was developed on three simulated processes of increasing
difficulty, all with leaks as the fault type. Scores are comparable within a
process but not between processes, since the scenarios and weights differ.
Table~\ref{tab:summary} summarizes the results. All numbers in this chapter
are re-evaluated with the final scoring rules of each process.

\begin{table}[tb]
  \centering
  \caption{Average score $J$ (lower is better) with PID control only and with
    the best generated supervisor. For the quadruple-tank process the best
    supervisor equals the perfect-detection oracle.}
  \label{tab:summary}
  \begin{tabular}{@{}lrrrr@{}}
    \toprule
    & \multicolumn{2}{c}{Development} & \multicolumn{2}{c}{Held-out} \\
    \cmidrule(lr){2-3}\cmidrule(l){4-5}
    Process & PID only & Supervisor & PID only & Supervisor \\
    \midrule
    Single tank      & 31\,133 & 7\,551 & 30\,945 & 10\,487 \\
    Two-tank cascade & 30\,775 & 8\,720 & 35\,274 & 13\,744 \\
    Quadruple tank   & 64\,865 & 2\,635 & 69\,569 &  2\,609 \\
    \bottomrule
  \end{tabular}
\end{table}

\section{Single tank}
\label{sec:single}
\paragraph{Setup.} One tank with a pump controlled by a PID loop
($K_p = 2$, $K_i = 0.5$, $K_d = 0.1$) at a sample time of 0.1~s, and an
unmeasured leak as the fault. The supervisor decides every 3~s on the last
30 samples of a 30~s episode. The development battery has eight scenarios:
fault-free, a standard leak, an early severe leak, a late mild leak, two
leaks that start and stop, and a fault-free and a leaking scenario with
sensor noise. The held-out battery has the same types with other parameters.

\paragraph{Results.} Of 62 trials, 16 were promoted. The final supervisor
reduced the development score by 76\% and the held-out score by 66\%
compared with PID alone. Its remaining cost is dominated by false positives:
across the development battery it has 192 false-positive samples against 34
missed ones, and one of its decisions in the fault-free scenario is a false
alarm.

\paragraph{Findings.} Early supervisors lowered the setpoint during a fault
and never raised it again, which led to the restoration term $R_s$ in
\eqref{eq:score}; the false-positive weight was also tuned during this
example. The held-out score was 39\% worse than the development score, a
first sign that the search fits the timing of the visible scenarios.

\section{Two-tank cascade}
\label{sec:two}
\paragraph{Setup.} Tank~1 drains by gravity into tank~2, and each tank has
its own pump and PID loop with the same gains as above. A leak in tank~1
alone moves the level of tank~2 by more than 2~m, so the supervisor must tell
propagated disturbances from genuine faults. The development battery has
eight scenarios: fault-free, a leak in either tank, simultaneous leaks,
severe persistent leaks in either tank, and a fault-free and a leaking
scenario with sensor noise.

\paragraph{Results.} Of 58 trials, 7 were promoted and 10 were rejected by
the regression guard: each improved the average score while making one
scenario much worse. The final supervisor reduced the development score by
72\% and the held-out score by 61\%. The held-out score was 58\% worse than
the development score.

\paragraph{Findings.} This example exposed specification gaming: a candidate
improved the average score by flagging anomalies almost constantly. A false
positive cost the same in the fault-free scenario as anywhere else, and the
loss there was outweighed by fewer missed anomalies elsewhere. This led to the
asymmetric false-positive cost and the per-scenario regression guard in
Section~\ref{sec:scoring}. The most instructive lesson recorded by the LLM was
that the supervisor's own mitigation corrupted its detection statistic.
Lowering the setpoint of a leaking tank lowers the pump effort needed to hold
it, so an effort-based detector saw the fault disappear, cleared the flag,
restored the setpoint and then detected the fault again. The resulting
oscillation showed up as matched bursts of missed anomalies and false
positives in every single-fault scenario. The LLM removed it by normalizing
the pump effort with the tank's own level, which makes the statistic
insensitive to setpoint changes. It also learned to separate propagated from
genuine downstream faults with a residual test: the excess effort in tank~2
minus a coupling gain times the excess effort in tank~1. Simultaneous faults
remained unsolved: every candidate that improved them lost robustness to
sensor noise, and consecutive trials often repeated nearly identical
proposals.

\section{Quadruple-tank process}
\label{sec:four}
\paragraph{Setup.} The quadruple-tank process \cite{johansson2000} from the
PC-Gym benchmark suite \cite{bloor2024pcgym} has four tanks and two pumps.
Each pump feeds one lower tank directly and, through a split valve, the upper
tank above the other lower tank. With the split fractions
$\gamma_1 = \gamma_2 = 0.2$, most of each pump's flow takes the delayed path,
which gives a non-minimum-phase process. Only the two lower levels are
measured. The model is integrated with the explicit Euler method at a 1~s
step over 600~s episodes. A leak multiplies the outlet area of a lower tank
by 1.8--4. There are six development and six held-out scenarios, and the
supervisor decides every 10~s on the last 50~s of telemetry.

\subsection{A plateau caused by the test bed}
The first training runs stalled at $J \approx 36\,300$. Analysis showed
three defects in the test bed rather than a limit of the search:
\begin{enumerate}
\item \emph{Faults were never applied.} They were injected through PC-Gym's
  environment wrapper, which silently integrated an unmodified copy of the
  model, so the first twenty or so trials ran without any faults. The model
  is now integrated directly.
\item \emph{Decision interval.} The supervisor decided only every 50~s, the
  same as its window length, and every development fault started on that
  grid. Each fault event therefore cost at least 51 missed and 51
  false-positive samples, however good the detector was. With decisions
  every 10~s on the same 50~s window, the unchanged champion's score fell
  from 36\,300 to 26\,323.
\item \emph{Loop pairing.} The relative gain of the diagonal pairing is
  \begin{equation*}
    \lambda = \frac{\gamma_1\gamma_2}{\gamma_1+\gamma_2-1} = -0.07 ,
  \end{equation*}
  so the loops should be paired off-diagonally \cite{johansson2000}. With the
  diagonal pairing, a leak in tank~1 saturated pump~1, whose flow mostly
  reaches tank~2; loop~2 then reduced pump~2 and drained tank~3, which is
  tank~1's main supply. The resulting 171 safety violations in the severe
  tank-1 scenario made up 39\% of the average score, and a grid search showed
  that no setpoint policy could remove them. After re-pairing and retuning
  the loops as PI controllers ($K_p = 40$, $K_i = 0.3$), PID control alone
  gave no violations in any scenario.
\end{enumerate}
At the plateau, the champion was within about 7\% of the best score any
supervisor could reach on the flawed test bed. Computing such a bound before
interpreting a plateau would have revealed the problem much earlier.

\subsection{Results after the corrections}
Table~\ref{tab:four} shows the corrected test bed. The best supervisors flag
every leak at the first decision after it starts and clear the flag at the
first decision after it stops, which gives exactly the oracle's number of
missed and false-positive samples in all twelve scenarios. Three insights
appear in their code and lessons: the loop regulating a leaking tank
saturates at the 12~V pump limit, and this is the one signature that the
supervisor's own setpoint reduction cannot hide; a falling level together
with an error above the healthy maximum reveals a leak several seconds
before saturation; and the rate at which a level rises separates the slow
creep of a leaking tank from refilling after a leak has stopped.

\begin{table}[tb]
  \centering
  \caption{Quadruple-tank process after the test-bed corrections.}
  \label{tab:four}
  \begin{tabular}{@{}lrr@{}}
    \toprule
    Supervisor & Development & Held-out \\
    \midrule
    PID only (and the seed supervisor)          & 64\,865 & 69\,569 \\
    Plateau champion, not retrained             & 10\,777 & 14\,941 \\
    Best generated (low reasoning, generation 2) &  2\,635 &  2\,609 \\
    Perfect-detection oracle                    &  2\,635 &  2\,609 \\
    \bottomrule
  \end{tabular}
\end{table}

\subsection{Effect of LLM reasoning}
The same search was run from the seed supervisor with three reasoning
settings (Figure~\ref{fig:reasoning} and Table~\ref{tab:reasoning}). With
high and low reasoning the search matched the oracle on both batteries within
three and two generations, respectively, and the held-out scores followed the
development scores. Without reasoning the search improved quickly at first
but stalled at 11\,837 from the third generation on, about 4.5 times the
oracle score. Three weaknesses stood out without reasoning: it missed 49--83
samples in each scenario with a mild leak, against 4--12 with reasoning; its lessons contained
incorrect physics, for example that a leak lowers the other loop's effort
through a shared outflow path that does not exist; and three of twelve
candidates in generations 4--6 called a built-in function that the sandbox
does not provide, after the prompt's list of allowed functions ended in
``etc.''. The list is now generated from the sandbox itself. With high
reasoning, 8 of 15 requests returned an empty response that was still
billed.

\begin{figure}[tb]
  \centering
  \includegraphics[width=\textwidth]{four_tank_reasoning_comparison}
  \caption{Score of the champion after each generation for three reasoning
    settings, on the development and held-out batteries (log scale). Dots
    are individual candidates; the three near $7\cdot 10^5$ crashed on every
    call. Generation~0 is the seed supervisor. The high-reasoning run used
    low reasoning in its third generation and an earlier, longer prompt in
    its first two.}
  \label{fig:reasoning}
\end{figure}

\begin{table}[tb]
  \centering
  \caption{Cost of reaching the oracle with different reasoning settings
    (DeepSeek off-peak prices, US dollars).}
  \label{tab:reasoning}
  \begin{tabular}{@{}lccc@{}}
    \toprule
    Reasoning & Generations & Output tokens & Cost \\
              & to oracle   & per request   & to oracle \\
    \midrule
    High & 3 & $\approx 61\,000$ & $\approx 0.63$ \\
    Low  & 2 & $\approx 23\,000$ & $\approx 0.11$ \\
    None & not reached in 6 & $\approx 2\,000$ & -- \\
    \bottomrule
  \end{tabular}
\end{table}

Each setting was run only once, and the runs differ slightly in their
prompts as noted in Figure~\ref{fig:reasoning}. The results therefore support
the conclusion that reasoning is needed and that low reasoning is enough for
this task, but not a precise ranking of low against high reasoning.
\TODO{Repeat each setting two or three times with the final prompt.}
